\section{Distillation Results: SFT es la base, DPO es el incremento}

Después de la human-data ablation, la clase regresa a Zephyr. Aquí hay que separar la leaderboard claim de la training ablation: la primera muestra la competitividad bajo un protocolo específico; solo la segunda responde de forma más directa cuál es la contribución de dSFT, chosen-only SFT y dDPO por separado.

\subsection{Lecture-time Leaderboard Evidence}

Zephyr-7B se basa en Mistral-7B; usa UltraChat para dSFT y después los pares chosen/rejected de UltraFeedback para dDPO. En clase se mostró que, en MT-Bench y AlpacaEval de ese momento, se acercaba a modelos más grandes o los superaba; estos resultados dependen de GPT-4 family judges, que son justamente lo que audita la siguiente sección.

\lecturefigure{slide-56.jpg}{Transición de sección hacia los distillation results}{Diapositiva oficial p.56}

\begin{knowledgebox}{Lectura de la figura: de la human-curated ablation a la AI-distilled recipe}
Aunque esta diapositiva contiene poca información, marca un cambio en la fuente de la evidencia: los modelos siguientes ya no aprenden principalmente de vendor-written demonstrations, sino de AI-generated dialogues/preferences. Al comparar, hay que tener en cuenta también la provenance.
\end{knowledgebox}

\lecturefigure{slide-57.jpg}{Resultados de Zephyr-7B en MT-Bench y AlpacaEval mostrados en clase}{Diapositiva oficial p.57}

\begin{knowledgebox}{Lectura de la figura: el recuadro rojo indica competitividad dentro del protocolo, no una victoria en todas las dimensiones}
La tabla incluye a la vez el tamaño en parámetros, MT-Bench y AlpacaEval. Primero se comparan los open models de la misma escala y luego el judge score frente a los proprietary systems; la figura no mide directamente la seguridad, la verificación de hechos, la retention de usuarios reales ni la robustness fuera de distribución.
\end{knowledgebox}

\subsection{Ablation: DPO-only, SFT-only y SFT+DPO}

La ablation aporta más información sobre el mecanismo. Aplicar DPO directamente sobre el base model da resultados muy pobres; solo con dSFT ya se obtiene la mayor parte de la ganancia; añadir dDPO después de dSFT mejora un poco más. Usar las chosen responses como SFT targets ordinarios tampoco sustituye por completo al preference objective.

\lecturefigure{slide-58.jpg}{Ablation de dSFT/dDPO en Zephyr}{Diapositiva oficial p.58}

\begin{knowledgebox}{Lectura de la figura: primero identificar las condiciones necesarias, después comparar los incrementos}
La fila DPO-only muestra que la preference optimization requiere una policy que ya tenga instruction-following; la fila SFT-only alcanza aproximadamente el 80--90\% del desempeño final; SFT+DPO es la mejor, pero con un incremento pequeño. La tabla respalda «DPO es refinement», no «DPO es inútil» ni «cualquier SFT basta».
\end{knowledgebox}

\teachervoice{El ponente dice explícitamente que DPO without SFT ``doesn't work at all'', que SFT ya aporta el 80--90\% del desempeño global y que DPO proporciona la ganancia restante. La chosen-only SFT tampoco equivale a aprovechar la pairwise preference.}

\subsection{Resumen de la sección}

Zephyr demuestra que los AI-generated data pueden sostener un open model sólido, pero también introducen en el pipeline el teacher/judge bias. La conclusión duradera de la ablation es: primero, SFT establece la base de instruction-following; después, la preference optimization realiza un refinamiento relativo.
